\documentclass[12pt]{article}
\usepackage{array} % Provides advanced control over table column formatting
\usepackage{xcolor}
\usepackage[margin=0.3in]{geometry}
\usepackage{bm}
\usepackage{scrextend}
\usepackage{titlesec}
\usepackage{amsmath}
\usepackage{booktabs}    % For professional table lines (\toprule, \midrule, \bottomrule)
\usepackage{subcaption}  % Required for the subtable environment

\let\frac\dfrac
\usepackage{cellspace}
\setlength{\cellspacetoplimit}{6pt}    % Minimum space above text
\setlength{\cellspacebottomlimit}{6pt} % minimum space belowe text

\changefontsizes{12pt}

\newcommand{\sectionbreak}{\clearpage}

\title{Community weighted traits}
\author{Prabu Ravindran}
\date{}

\begin{document}
	\maketitle
	
	In this document, we will demonstrate one-level and two-level community weighting with a single plot. 
			
	\begin{table}[htbp]
		\centering
		\caption{Plots and coverages.}
		\label{tab:main_grid}
		
		% Top-Left Table
		\begin{subtable}[b]{0.48\textwidth}
			\centering
			\caption{Plot for \textbf{one-level} weighting.}
			\label{tab:table_a}
			\begin{tabular}{|c Sc|}
				\hline 
				$A_1^{ss_1}$  & $\mathcolor{red}{\bm{A_2^{ss_1}}}$  \\ 
				$B_1^{ss_1}$  & $C_1^{ss_1}$  \\ \hline
			\end{tabular}
		\end{subtable}
		\hfill % Adds horizontal space between the two columns
		% Top-Right Table
		\begin{subtable}[b]{0.48\textwidth}
			\centering
			\caption{Plot for \textbf{two-level} weighting.}
			\label{tab:table_b}
			\begin{tabular}{|c c||c Sc|}
				\hline 
				$A_1^{ss_1}$  & $\mathcolor{red}{\bm{A_2^{ss_1}}}$ & $A_1^{ss_2}$ & $C_1^{ss_2}$ \\ 
				$B_1^{ss_1}$  & $C_1^{ss_1}$ & $D_1^{ss_2}$ & $\mathcolor{blue}{\bm{D_2^{ss_2}}}$ \\ \hline
			\end{tabular}
		\end{subtable}
		
		\vspace{0.5cm} % Adds vertical spacing between the two rows
		
		% --- SECOND ROW ---
		% Bottom-Left Table
		\begin{subtable}[b]{0.48\textwidth}
			\centering
			\caption{Coverage for \textbf{one-level} weighting}
			\label{tab:table_c}
			\begin{tabular}{c c Sc}
				\toprule 
				\texttt{SubsampleID} & \texttt{CoverID} & \texttt{Coverage} \\
				\midrule 
				$ss_{1}$ & $A$  & $w_{A}^{ss_1}$  \\
				$ss_{1}$ & $B$  & $w_{B}^{ss_1}$  \\
				$ss_{1}$ & $C$  & $w_{C}^{ss_1}$  \\
				\bottomrule
			\end{tabular}
%			\begin{tabular}{c c c c c Sc Sc Sc}
%				\toprule
%				\texttt{SubsampleID} & \texttt{SubsampleWeight} & \texttt{CoverID} & \texttt{CoverWeight} & \texttt{TraitMean} & \texttt{TraitSdev}\\
%				\midrule  
%				$ss_1$ & $W^{ss_1}$ & $A_1$ & $w_{A}^{ss_1}$ & $\mu_{A_1}^{ss_1}$ & $\sigma_{A_1}^{ss_1}$ \\
%				$ss_1$ & $W^{ss_1}$ & $A_2$ & $w_{A}^{ss_1}$ & $\mathcolor{red}{\bm{\mu_{A_2}^{ss_1}}}$ & $\mathcolor{red}{\bm{\sigma_{A_2}^{ss_1}}}$ \\ 
%				$ss_1$ & $W^{ss_1}$ & $B_1$ & $w_{B}^{ss_1}$ & $\mu_{B_1}^{ss_1}$ & $\sigma_{B_1}^{ss_1}$ \\
%				$ss_1$ & $W^{ss_1}$ & $C_1$ & $w_{C}^{ss_1}$ & $\mu_{C_1}^{ss_1}$ & $\sigma_{C_1}^{ss_1}$ \\
%				\bottomrule
%			\end{tabular}
		\end{subtable}
		\hfill % Adds horizontal space
		% Bottom-Right Table
		\begin{subtable}[b]{0.48\textwidth}
			\centering
			\caption{Coverage for \textbf{two-level} weighting}
			\label{tab:table_d}
			\begin{tabular}{c c Sc}
				\toprule 
				\texttt{SubsampleID} & \texttt{CoverID} & \texttt{Coverage} \\ 
				\midrule  
				$ss_{1}$ & $A$  & $w_{A}^{ss_1}$  \\ 
				$ss_{1}$ & $B$  & $w_{B}^{ss_1}$  \\
				$ss_{1}$ & $C$  & $w_{C}^{ss_1}$  \\ 
				\hline
				$ss_{2}$ & $A$  & $w_{A}^{ss_2}$  \\ 
				$ss_{2}$ & $C$  & $w_{C}^{ss_2}$  \\
				$ss_{2}$ & $D$  & $w_{D}^{ss_2}$  \\ 
				\bottomrule
			\end{tabular}
		\end{subtable}
		
	\end{table}
		
	
		
	\section*{One level weighting}
	% The plot
	
	A cartoon of a plot and the measured \texttt{CoverID} are shown in~\ref{tab:l1_plot}. There are $3$ \texttt{CoverID}s recorded in the plot. \texttt{CoverID}$={A}$ has two ``individuals"; while \texttt{CoverID}s $B, C$ have one ``individual" each. 
	
	\begin{table}[htbp]
		\centering 
		\caption{The plot}
		\label{tab:l1_plot} 
		\begin{tabular}{|c|Sc|}
			\hline 
			$A_1^{ss_1}$  & $\mathcolor{red}{\bm{A_2^{ss_1}}}$  \\ \hline
			$B_1^{ss_1}$  & $C_1^{ss_1}$  \\ \hline
		\end{tabular}
	\end{table}
	
	% The plot composition
	Field notes for the above plot will record the cover information for each \texttt{CoverID}. This can be recorded as \texttt{CoverPercent} or \texttt{CoverClass}. For \texttt{CoverClass} specification, coverage bins must also be provided. The column \texttt{SubsampleID} 
	\begin{table}[htbp]
		\centering 
		\caption{Field notes for plot composition}
		\label{tab:l1_field_notes_percent} 
		\begin{tabular}{|c|c|Sc|}
			\hline 
			\texttt{SubsampleID} & \texttt{CoverID} & \texttt{CoverPercent} \\ \hline 
			$ss_{1}$ & $A$  & $w_{A}^{ss_1}=40\%$  \\ \hline
			$ss_{1}$ & $B$  & $w_{B}^{ss_1}=20\%$  \\ \hline
			$ss_{1}$ & $C$  & $w_{C}^{ss_1}=10\%$  \\ \hline
		\end{tabular}
	\end{table} 
	
	\clearpage
	 
	% The DataFrame
	\begin{table}[htbp]
		\centering 
		\caption{The dataframe. \texttt{Trait} standard deviation column is not shown.}
		\label{tab:l1_dataframe} 
		\begin{tabular}{|c|c|c|c|c|Sc|Sc|Sc|}
			\hline 
			\texttt{SampleID} & \texttt{SubsampleID} & \texttt{SubsampleWeight} & \texttt{CoverID} & \texttt{CoverWeight} & \texttt{TraitMean} & \texttt{TraitSdev}\\ \hline 
			$S_1$  & $ss_1$ & $W^{ss_1}$ & $A_1$ & $w_{A}^{ss_1}=40$ & $\mu_{A_1}^{ss_1}$ & $\sigma_{A_1}^{ss_1}$ \\ \hline
			$S_1$  & $ss_1$ & $W^{ss_1}$ & $A_2$ & $w_{A}^{ss_1}=40$ & $\mathcolor{red}{\bm{\mu_{A_2}^{ss_1}}}$ & $\mathcolor{red}{\bm{\sigma_{A_2}^{ss_1}}}$ \\ \hline
			$S_1$  & $ss_1$ & $W^{ss_1}$ & $B_1$ & $w_{B}^{ss_1}=20$ & $\mu_{B_1}^{ss_1}$ & $\sigma_{B_1}^{ss_1}$ \\ \hline
			$S_1$  & $ss_1$ & $W^{ss_1}$ & $C_1$ & $w_{C}^{ss_1}=10$ & $\mu_{C_1}^{ss_1}$ & $\sigma_{C_1}^{ss_1}$ \\ \hline
		\end{tabular}
	\end{table}
		
	% Weighted traits for subsample
	\begin{table}[htbp]
		\centering 
		\caption{Algebra for going from dataframe to \texttt{SubsampleTrait}}
		\label{tab:l1_subsample_trait}
		\begin{tabular}{|c|c|c|Sc|}
			\hline 
			\texttt{SampleID} & \texttt{SubsampleID} & \texttt{SubsampleWeight} & \texttt{SubsampleTrait} \\ \hline 
			$S_1$  & $ss_1$ & $W^{ss_1}$ & $t^{ss_1}(t_{A_1}^{ss_1}, t_{B_1}^{ss_1}, t_{C_1}^{ss_1})$ = $\frac{w_{A}^{ss_1}t_{A_1}^{ss_1} + w_{B}^{ss_1}t_{B_1}^{ss_1} + w_{C}^{ss_1}t_{C_1}^{ss_1}}{w_{A}^{ss_1} + w_{B}^{ss_1} + w_{C}^{ss_1}}$  \\ \hline
			$S_1$  & $ss_1$ & $W^{ss_1}$ & $t^{ss_1}(\mathcolor{red}{\bm{t_{A_2}^{ss_1}}}, t_{B_1}^{ss_1}, t_{C_1}^{ss_1}) = \frac{w_{A}^{ss_1}\mathcolor{red}{t_{A_2}^{ss_1}} + w_{B}^{ss_1}t_{B_1}^{ss_1} + w_{C}^{ss_1}t_{C_1}^{ss_1}}{w_{A}^{ss_1} + w_{B}^{ss_1} + w_{C}^{ss_1}}$ \\ \hline
		\end{tabular}
	\end{table}

	% Weighted traits for sample
	\begin{table}[htbp]
		\centering
		\caption{Algebra for going from \texttt{SubsampleTrait} to \texttt{CWTrait}}
		\label{tab:l1_sample_trait}
		\begin{tabular}{|c|Sc|}
			\hline 
			\texttt{SampleID} & \texttt{CWTrait}\\ \hline 
			$S_1$  & $T = \frac{W^{ss_1}t^{ss_1}(t_{A_1}^{ss_1}, t_{B_1}^{ss_1}, t_{C_1}^{ss_1})}{W^{ss_1}}$ \\ \hline
			$S_1$  & $T = \frac{W^{ss_1}t^{ss_1}(\mathcolor{red}{\bm{t_{A_2}^{ss_1}}}, t_{B_1}^{ss_1}, t_{C_1}^{ss_1})}{W^{ss_1}}$ \\ \hline
		\end{tabular}
	\end{table}	
	
	\pagebreak 
	
	\section*{Notes} 
	\begin{itemize}
		\item
		\textbf{Trait uncertainty}: The value used for (say) $t_{A_1}^{ss_1}$ can be sampled from a distribution specified with columns for mean trait value and the associated standard deviation. In our context, the mean and deviations are typically produced from leaf-level model predictions. The trait standard deviation column is omitted in the tables for clarity. 
		
		\item 
		\textbf{Cover replicate uncertainty}: The two trait values for \texttt{CWTrait} in table~\ref{tab:l1_sample_trait} (and column \texttt{SubsampleTrait} in table~\ref{tab:l1_subsample_trait}) are due to choice of the replicate for \texttt{CoverID} $A$. This is ``specified" by the number of replicates for each \texttt{CoverID}. 
		
		\item
		\textbf{Cover weight uncertainty}: Ordinal cover classes and percent cover are two approaches for specifying \texttt{Coverweight} (e.g. $w_{A}^{ss_1}$). Ordinal cover class specification requires binning the range $[0, 100]$ and the bin membership for the considered \texttt{CoverID}. For percent cover, \texttt{CoverID} is specified as a percentage (range: $[0, 100]$) or a proportion (range: $[0, 1]$). Typically, human visual estimates are used for the \texttt{CoverWeight}, and the imprecision of human observation leads to uncertainties in \texttt{CoverWeight}.   
		
%		\item 
%		\textbf{Sampling \texttt{CWTrait}}: Propagating the \texttt{CWTrait} uncertainty is critical for robust model training and evaluation. Specifically, \texttt{CWTrait} ($\widetilde{y}$) can be computed using sampled values for $t_{c}^{s}$ and $w_{c}^{s}$, where $s$ is \texttt{SubsampleID} (say $ss_1$) and $c$ is the \texttt{CoverID} (e.g., $A$ for $w$ and $A_1$ for $t$). The spectra-trait pair $(X, \widetilde{y})$ is used for model training, with different sampled values of $\widetilde{y}$ for a \texttt{SubsampleID} during various folds of cross-validation.      
		
		\item 
		\textbf{Montecarlo \texttt{CWTrait} estimation}:
		Monte-Carlo estimation provides a general, tractable estimate of \texttt{CWTrait} (${y}_{\mu}^{cw}$) (and its standard deviation, $y_{\sigma}^{cw}$). Given a dataset and distributions for $t_c^{s}$ and $w_c^{s}$ (see below), we can compute $\{\widetilde{y_1}, \ldots \widetilde{y_M}\}$ ($\widetilde{y_i}$ uses samples  $(\widetilde{(t_c^{s})_i}$ $\widetilde{(w_c^{s})_i})$ in the formulas in the tables). The $\{\widetilde{y_i}\}$ are reduced to yield $y_{\mu}^{cw}$ and $y_{\sigma}^{cw}$. The $y_{\mu}^{cw}$ is used as the ``true" value for model evaluation and in the predicted-vs-true plot; $y_{\sigma}^{cw}$ provides the error bar for the ``true" value in a predicted-vs-true plot; and sample with mean $y_{\mu}^{cw}$ and standard deviation $y_{\sigma}^{cw}$ is used during model training.
		
		\item 
		\textbf{Ordinal \texttt{CoverWeight} specification}: 
		\begin{itemize}
			\item
			Bin index (integer) $i \in \{ 1, \ldots B\}$.
			\item 
			The end points for the $B$ bins $e_0, e_1, \ldots, e_B$, $e_b \in [0, 100]$; bin mode offsets, $o_1, \ldots o_B \in [0, 1]$; and bin concentrations for ``peakiness", $k_1, \ldots, k_B$ (table~\ref{tab:ordinal_cover}). The ``representative" value for bin $b$ is $r_b = e_{b - 1} + o_b(e_b - e_{b - 1})$. 
			\begin{table}[htbp]
				\centering
				\caption{Specification for ordinal cover}
				\label{tab:ordinal_cover}
				\begin{tabular}{|c|c|c|c|Sc|}
					\hline 
					\texttt{BinIndex} & \texttt{LeftEnd} & \texttt{RightEnd} & \texttt{Offset} & \texttt{Concentration}\\ \hline 
					$1$  & $e_0$ & $e_1$ & $o_1$ & $k_1$ \\ \hline
					$\cdots$ & $\cdots$ & $\cdots$ & $\cdots$ & $\cdots$ \\ \hline
					$B$  & $e_{B-1}$ & $e_B$ & $o_B$ & $k_B$ \\ \hline
				\end{tabular}
			\end{table}	
			\item
			A mode $r_b$, concentration $k_b$, support $[\texttt{max}(0, r_{b - 1}), \texttt{min}(r_{b + 1}, 100)]$, \texttt{Beta} distribution is used for sampling $\widetilde{(w_c^{s})_i}$, if \texttt{CoverWeight} = $b$. 
		\end{itemize} 
					
				
					
		\item
		\textbf{Percent \texttt{CoverWeight} specification}:  
		\begin{itemize}
			\item
			Percentage (integer) $p \in [0, 100]$. 
			\item 
			The end points for $B$ bins: $e_0, e_1, \ldots, e_B \in [0, 100]$; and uncertainties $u_1, \ldots u_B \in [0, 100]$. $u_b$ is the uncertainty, if $p \in [e_{b - 1}, e_b)$ and corresponds to one standard deviation. 
			\begin{table}[htbp]
				\centering
				\caption{Specification for percent cover}
				\label{tab:percent_cover}
				\begin{tabular}{|c|c|c|Sc|}
					\hline 
					\texttt{BinIndex} & \texttt{LeftEnd} & \texttt{RightEnd} & \texttt{Uncertainty} \\ \hline 
					$1$  & $e_0$ & $e_1$ & $u_1$ \\ \hline
					$\cdots$ & $\cdots$ & $\cdots$ & $\cdots$ \\ \hline
					$B$  & $e_{B-1}$ & $e_B$ & $u_B$ \\ \hline
				\end{tabular}
			\end{table}
			\item 
			A \texttt{TruncNormal} with mode $p$ and range $[\texttt{max}(0, p - u_b), \texttt{min}(p + u_b, 100)]$ is used for sampling $\widetilde{(w_c^{s})_i}$, if $p \in [e_{b - 1}, e_{b})$.
			
		\end{itemize} 
			
	\end{itemize}
	
	%%%%%%%%%%%%%%%%%%%%%%%%%%%%%%%%%%%%%%%%%%%%%%%%%%%%%%%%%%%%%%%%%%%%%%%%%
	
	\pagebreak
	
	\section*{Two level weighting}
	
	% The plot
	\begin{table}[htbp]
		\centering 
		\caption{The plot}
		\label{tab:l2_plot}  
		\begin{tabular}{|c|c||c|Sc|}
			\hline 
			$A_1^{ss_1}$  & $\mathcolor{red}{\bm{A_2^{ss_1}}}$ & $A_1^{ss_2}$ & $C_1^{ss_2}$ \\ \hline
			$B_1^{ss_1}$  & $C_1^{ss_1}$ & $D_1^{ss_2}$ & $\mathcolor{blue}{\bm{D_2^{ss_2}}}$ \\ \hline
		\end{tabular}
	\end{table}
	
	% The plot composition
	\begin{table}[htbp]
		\centering  
		\caption{Field notes for plot composition}
		\label{tab:l2_field_notes}
		\begin{tabular}{|c|c|Sc|}
			\hline 
			\texttt{SubsampleID} & \texttt{CoverID} & \texttt{CoverPercent} \\ \hline 
			$ss_{1}$ & $A$  & $w_{A}^{ss_1}$  \\ \hline
			$ss_{1}$ & $B$  & $w_{B}^{ss_1}$  \\ \hline
			$ss_{1}$ & $C$  & $w_{C}^{ss_1}$  \\ \hline
			
			$ss_{2}$ & $A$  & $w_{A}^{ss_2}$  \\ \hline
			$ss_{2}$ & $C$  & $w_{C}^{ss_2}$  \\ \hline
			$ss_{2}$ & $D$  & $w_{D}^{ss_2}$  \\ \hline
		\end{tabular}
	\end{table} 
	
	\pagebreak
	
	
	% The DataFrame
	\begin{table}[htbp]
		\centering 
		\caption{The dataframe. \texttt{Trait} standard deviation column is not shown.}
		\label{tab:l2_dataframe}
		\begin{tabular}{|c|c|c|c|c|Sc|}
			\hline 
			\texttt{SampleID} & \texttt{SubsampleID} & \texttt{SubsampleWeight} & \texttt{CoverID} & \texttt{CoverWeight} & \texttt{Trait}\\ \hline 
			$S_1$  & $ss_1$ & $W^{ss_1}$ & $A_1$ & $w_{A}^{ss_1}$ & $t_{A_1}^{ss_1}$ \\ \hline
			$S_1$  & $ss_1$ & $W^{ss_1}$ & $A_2$ & $w_{A}^{ss_1}$ & $\mathcolor{red}{\bm{t_{A_2}^{ss_1}}}$ \\ \hline
			$S_1$  & $ss_1$ & $W^{ss_1}$ & $B_1$ & $w_{B}^{ss_1}$ & $t_{B_1}^{ss_1}$ \\ \hline
			$S_1$  & $ss_1$ & $W^{ss_1}$ & $C_1$ & $w_{C}^{ss_1}$ & $t_{C_1}^{ss_1}$ \\ \hline
		
			$S_1$  & $ss_2$ & $W^{ss_2}$ & $A_1$ & $w_{A}^{ss_2}$ & $t_{A_1}^{ss_2}$ \\ \hline
			$S_1$  & $ss_2$ & $W^{ss_2}$ & $C_1$ & $w_{C}^{ss_2}$ & $t_{C_1}^{ss_2}$ \\ \hline
			$S_1$  & $ss_2$ & $W^{ss_2}$ & $D_1$ & $w_{D}^{ss_2}$ & $t_{D_1}^{ss_2}$ \\ \hline
			$S_1$  & $ss_2$ & $W^{ss_2}$ & $D_2$ & $w_{D}^{ss_2}$ & $\mathcolor{blue}{\bm{t_{D_2}^{ss_2}}}$ \\ \hline
		\end{tabular}
	\end{table}
	
	% Weighted traits for subsample
	\begin{table}[htbp]
		\centering 
		\caption{Algebra for going from dataframe to \texttt{SubsampleTrait}}
		\label{tab:l2_subsample_trait}
		\begin{tabular}{|c|c|c|Sc|}
			\hline 
			\texttt{SampleID} & \texttt{SubsampleID} & \texttt{SubsampleWeight} & \texttt{SubsampleTrait} \\ \hline 
			
			$S_1$  & $ss_1$ & $W^{ss_1}$ & $t^{ss_1}(t_{A_1}^{ss_1}, t_{B_1}^{ss_1}, t_{C_1}^{ss_1}) = \frac{w_{A}^{ss_1}t_{A_1}^{ss_1} + w_{B}^{ss_1}t_{B_1}^{ss_1} + w_{C}^{ss_1}t_{C_1}^{ss_1}}{w_{A}^{ss_1} + w_{B}^{ss_1} + w_{C}^{ss_1}}$ \\ \hline
			$S_1$  & $ss_1$ & $W^{ss_1}$ & $t^{ss_1}(\mathcolor{red}{\bm{t_{A_2}^{ss_1}}}, t_{B_1}^{ss_1}, t_{C_1}^{ss_1}) = \frac{w_{A}^{ss_1}\mathcolor{red}{\bm{t_{A_2}^{ss_1}}} + w_{B}^{ss_1}t_{B_1}^{ss_1} + w_{C}^{ss_1}t_{C_1}^{ss_1}}{w_{A}^{ss_1} + w_{B}^{ss_1} + w_{C}^{ss_1}}$ \\ \hline
			
			$S_1$  & $ss_2$ & $W^{ss_2}$ & $t^{ss_2}(t_{A_1}^{ss_2}, t_{C_1}^{ss_2}, t_{D_1}^{ss_2}) = \frac{w_{A}^{ss_2}t_{A_1}^{ss_2} + w_{C}^{ss_2}t_{C_1}^{ss_2} + w_{D}^{ss_2}t_{D_1}^{ss_2}}{w_{A}^{ss_2} + w_{C}^{ss_2} + w_{D}^{ss_2}}$ \\ \hline
			$S_1$  & $ss_2$ & $W^{ss_2}$ & $t^{ss_2}(t_{A_1}^{ss_2}, t_{C_1}^{ss_2}, \mathcolor{blue}{\bm{t_{D_2}^{ss_2}}}) = \frac{w_{A}^{ss_2}t_{A_1}^{ss_2} + w_{C}^{ss_2}t_{C_1}^{ss_2} + w_{D}^{ss_2}\mathcolor{blue}{\bm{t_{D_2}^{ss_2}}}}{w_{A}^{ss_2} + w_{C}^{ss_2} + w_{D}^{ss_2}}$ \\ \hline
			
			
		\end{tabular}
	\end{table}
	
	% Weighted traits for sample
	\begin{table}[htbp]
		\centering
		\caption{Algebra for going from \texttt{SubsampleTrait} to \texttt{SampleTrait}}
		\label{tab:l2_sample_trait}
		\begin{tabular}{|c|Sc|}
			\hline 
			\texttt{SampleID} & \texttt{CWTrait}\\ \hline 
			$S_1$  & $T = \frac{W^{ss_1}t^{ss_1}(t_{A_1}^{ss_1}, t_{B_1}^{ss_1}, t_{C_1}^{ss_1}) + W^{ss_2}t^{ss_2}(t_{A_1}^{ss_2}, t_{C_1}^{ss_2}, t_{D_1}^{ss_2})}{W^{ss_1} + W^{ss_2}}$ \\ \hline
			$S_1$  & $T = \frac{W^{ss_1}t^{ss_1}(\mathcolor{red}{\bm{t_{A_2}^{ss_1}}}, t_{B_1}^{ss_1}, t_{C_1}^{ss_1}) + W^{ss_2}t^{ss_2}(t_{A_1}^{ss_2}, t_{C_1}^{ss_2}, t_{D_1}^{ss_2})}{W^{ss_1} + W^{ss_2}}$ \\ \hline
			$S_1$  & $T = \frac{W^{ss_1}t^{ss_1}(t_{A_1}^{ss_1}, t_{B_1}^{ss_1}, t_{C_1}^{ss_1}) + W^{ss_2}t^{ss_2}(t_{A_1}^{ss_2}, t_{C_1}^{ss_2}, \mathcolor{blue}{\bm{t_{D_2}^{ss_2}}})}{W^{ss_1} + W^{ss_2}}$ \\ \hline
			$S_1$  & $T = \frac{W^{ss_1}t^{ss_1}(\mathcolor{red}{\bm{t_{A_2}^{ss_1}}}, t_{B_1}^{ss_1}, t_{C_1}^{ss_1}) + W^{ss_2}t^{ss_2}(t_{A_1}^{ss_2}, t_{C_1}^{ss_2}, \mathcolor{blue}{\bm{t_{D_2}^{ss_2}}})}{W^{ss_1} + W^{ss_2}}$ \\ \hline
		\end{tabular}
	\end{table}	
	
	
\end{document}
